\chapter{Isentropic Flow}

Isentropic flow describes compressible flow through a nozzle or diffuser where there is no heat transfer and no friction. The process is both adiabatic and reversible. It is the idealized model that connects the local Mach number at any cross-section to the stagnation (total) conditions of the flow. All flow properties at a given Mach number can be expressed as simple functions of $\gamma$ and $M$, making it a foundational tool in nozzle design and supersonic aerodynamics. This module was recreated from the legacy toolkit codebase originally developed by Gowtham Sivaraman \cite{sivaraman2015}.

\section{Nomenclature}

\begin{center}
\begin{tabular}{@{}cl@{}}
\toprule
Symbol & Description \\
\midrule
$\gamma$ & Specific heat ratio \\
$M$ & Mach number \\
$\frac{T}{T_0}$ & Static to Stagnation temperature ratio \\
$\frac{P}{P_0}$ & Static to Stagnation pressure ratio \\
$\frac{\rho}{\rho_0}$ & Static to Stagnation density ratio \\
$\frac{A}{A^*}$ & Local area to Throat area ratio \\
$\mu$ & Mach angle \\
$\nu$ & Prandtl--Meyer expansion angle \\
\bottomrule
\end{tabular}
\end{center}

\section{Governing Idea}

Every isentropic-flow property is a function of $\gamma$ and $M$ alone. The
solver therefore treats $M$ as the pivot variable: whichever quantity the
user supplies is first inverted to recover $M$, after which all remaining
properties are evaluated directly as explicit functions of $M$ and $\gamma$.
The problem is split into seven cases depending on which property is given.

\section{Calculation Cases}

Depending on the input parameter provided by the user, the isentropic flow properties are computed using one of the seven cases described below.

\subsection*{Case 1: Given $\gamma$ and $M$}

\textbf{Given:}
\begin{itemize}[nosep]
    \item Specific heat ratio $\gamma > 1$
    \item Mach number $M \ge 0$
\end{itemize}

\textbf{Calculation:} All properties follow directly from explicit analytical relations:
\begin{equation}
\frac{T}{T_0} = \frac{1}{1 + \frac{\gamma-1}{2}M^2}
\end{equation}

\begin{equation}
\frac{P}{P_0} = \left(1 + \frac{\gamma-1}{2}M^2\right)^{\frac{\gamma}{1-\gamma}}
\end{equation}

\begin{equation}
\frac{\rho}{\rho_0} = \left(1 + \frac{\gamma-1}{2}M^2\right)^{\frac{1}{1-\gamma}}
\end{equation}

\begin{equation}
\frac{A}{A^*} = \frac{1}{M}\left[\frac{1+\left(\frac{\gamma-1}{2}\right)M^2}{\frac{\gamma+1}{2}}\right]^{\frac{\gamma+1}{2(\gamma-1)}}
\end{equation}

For supersonic flow ($M \ge 1$), the Mach angle and Prandtl--Meyer expansion angle follow as:
\begin{equation}
\mu = \sin^{-1}\left(\frac{1}{M}\right)
\end{equation}

\begin{equation}
\nu = \left(\frac{\gamma+1}{\gamma-1}\right)^{\frac{1}{2}}\tan^{-1}\left[\frac{\gamma-1}{\gamma+1}(M^2-1)\right]^{\frac{1}{2}} - \tan^{-1}(M^2-1)^{\frac{1}{2}}
\end{equation}
For subsonic flow ($0 \le M < 1$), the values of $\mu$ and $\nu$ are mathematically undefined and not computed.


\subsection*{Case 2: Given $\gamma$ and $\frac{T}{T_0}$}

\textbf{Given:}
\begin{itemize}[nosep]
    \item Specific heat ratio $\gamma > 1$
    \item Stagnation temperature ratio $0 < \frac{T}{T_0} < 1$
\end{itemize}

\textbf{Calculation:} The Mach number is solved explicitly in closed form:
\begin{equation}
M = \sqrt{\frac{2}{\gamma-1}\left[\left(\frac{1}{T/T_0}\right)-1\right]}
\end{equation}


\subsection*{Case 3: Given $\gamma$ and $\frac{P}{P_0}$}

\textbf{Given:}
\begin{itemize}[nosep]
    \item Specific heat ratio $\gamma > 1$
    \item Stagnation pressure ratio $0 < \frac{P}{P_0} < 1$
\end{itemize}

\textbf{Calculation:} The Mach number is solved explicitly in closed form:
\begin{equation}
M = \sqrt{\frac{2}{\gamma-1}\left[\left(\frac{P}{P_0}\right)^{\frac{1-\gamma}{\gamma}}-1\right]}
\end{equation}


\subsection*{Case 4: Given $\gamma$ and $\frac{\rho}{\rho_0}$}

\textbf{Given:}
\begin{itemize}[nosep]
    \item Specific heat ratio $\gamma > 1$
    \item Stagnation density ratio $0 < \frac{\rho}{\rho_0} < 1$
\end{itemize}

\textbf{Calculation:} The Mach number is solved explicitly in closed form:
\begin{equation}
M = \sqrt{\frac{2 \times \left[\left(\frac{1}{\rho/\rho_0}\right)^{\gamma-1}-1\right]}{\gamma-1}}
\end{equation}


\subsection*{Case 5: Given $\gamma$ and $\frac{A}{A^*}$}

\textbf{Given:}
\begin{itemize}[nosep]
    \item Specific heat ratio $\gamma > 1$
    \item Area ratio $\frac{A}{A^*} \ge 1$ (values strictly below 1 are unphysical)
\end{itemize}

\textbf{Calculation:} No analytical closed-form inversion exists; the Mach number is solved iteratively using the Newton--Raphson method. Let the residual function and its derivative with respect to Mach number be defined as:
\begin{equation}
f(M) = \frac{1}{M}\left[\frac{1+\left(\frac{\gamma-1}{2}\right)M^2}{\frac{\gamma+1}{2}}\right]^{\frac{\gamma+1}{2(\gamma-1)}} - \frac{A}{A^*} = 0
\end{equation}
\begin{equation}
f'(M) = \frac{\left(\frac{2}{\gamma+1}\right)^{\frac{\gamma+1}{2(\gamma-1)}}\left[M^2\left(\frac{\gamma+1}{2}\right)\left(1+\frac{\gamma-1}{2}M^2\right)^{\frac{3-\gamma}{2(\gamma-1)}} - \left(1+\frac{\gamma-1}{2}M^2\right)^{\frac{\gamma+1}{2(\gamma-1)}}\right]}{M^2}
\end{equation}

Beginning with an initial guess $M_0$, successive estimates are calculated via:
\begin{equation}
M_{new} = M - \frac{f(M)}{f'(M)}
\end{equation}
The iteration continues until convergence, defined when $|M_{new} - M| < 10^{-6}$ (with a limit of $10{,}000$ iterations).

Due to the non-monotonic nature of the area ratio, two physical solutions exist:
\begin{itemize}[nosep]
    \item \textbf{Subsonic branch:} Solved using a seed value of $M_0 = 0.00001$.
    \item \textbf{Supersonic branch:} Solved using a seed value of $M_0 = 2.0$.
\end{itemize}


\subsection*{Case 6: Given $\gamma$ and $\mu$}

\textbf{Given:}
\begin{itemize}[nosep]
    \item Specific heat ratio $\gamma > 1$
    \item Mach angle $0^\circ < \mu \le 90^\circ$ (where $\mu$ in radians is $\mu_{rad} = \frac{\pi}{180}\mu_{deg}$)
\end{itemize}

\textbf{Calculation:} The Mach number is solved explicitly in closed form:
\begin{equation}
M = \frac{1}{\sin\mu_{rad}}
\end{equation}
Since $\mu \le 90^\circ$, this solution is strictly supersonic ($M \ge 1$).


\subsection*{Case 7: Given $\gamma$ and $\nu$}

\textbf{Given:}
\begin{itemize}[nosep]
    \item Specific heat ratio $\gamma > 1$
    \item Prandtl--Meyer expansion angle $0^\circ \le \nu < \nu_{max}$, where the maximum expansion angle $\nu_{max}$ is:
    \begin{equation}
    \nu_{max} = \left(\sqrt{\frac{\gamma+1}{\gamma-1}} - 1\right) \times 90^\circ
    \end{equation}
    Let $\nu_{rad} = \frac{\pi}{180}\nu_{deg}$ be the input in radians.
\end{itemize}

\textbf{Calculation:} No analytical closed-form inversion exists; the Mach number is solved iteratively using the Newton--Raphson method. Let the residual function and its derivative with respect to Mach number be defined as:
\begin{equation}
f(M) = \left(\frac{\gamma+1}{\gamma-1}\right)^{\frac{1}{2}}\tan^{-1}\left[\frac{\gamma-1}{\gamma+1}(M^2-1)\right]^{\frac{1}{2}} - \tan^{-1}(M^2-1)^{\frac{1}{2}} - \nu_{rad} = 0
\end{equation}
\begin{equation}
f'(M) = \frac{M\sqrt{M^2-1}}{1+\frac{\gamma-1}{2}M^2}
\end{equation}

Beginning with an initial guess of $M_0 = 2.0$ (or $M_0 = 1.05$ for small angles $\nu_{deg} < 1.0$), successive estimates are calculated via:
\begin{equation}
M_{new} = M - \frac{f(M)}{f'(M)}
\end{equation}
The iteration continues until convergence, defined when $|M_{new} - M| < 10^{-6}$. If a step results in $M_{new} \le 1$, the solver resets the trial to a midpoint $M_{new} = \frac{M + 1.0}{2.0}$ to ensure $M > 1$ remains true.




\section{Program Flow and Architecture}

In isentropic flow, $M$ is the single pivot variable, every input is first inverted to recover $M$, then all outputs are computed directly from $M$ and $\gamma$. The diagram below shows this convergent flow structure.

\begin{center}
\begin{tikzpicture}[x=1cm,y=1cm,
    inp/.style={rectangle, draw=blue!70!black, thick, fill=blue!6, align=center,
                text width=4.5cm, minimum height=0.65cm, font=\footnotesize\sffamily},
    eng/.style={rectangle, draw=orange!80!black, thick, fill=orange!8, align=center,
                text width=3.4cm, font=\small\sffamily},
    outn/.style={rectangle, draw=green!60!black, thick, fill=green!6, align=center,
                text width=2.2cm, minimum height=0.65cm, font=\footnotesize\sffamily},
    note/.style={rectangle, draw=gray!60, dashed, fill=gray!5, align=center,
                 text width=2.0cm, minimum height=0.55cm, font=\scriptsize\sffamily},
    arr/.style={-{Stealth[scale=1.0]}, thick, draw=black!70},
    darr/.style={-{Stealth[scale=1.0]}, thick, dashed, draw=black!50},
    bus/.style={thick, draw=orange!70!black}
]
%--- Inputs ---
\node[inp] (i0) at (0, 0.0) {Case 1: $M$ (direct)};
\node[inp] (i1) at (0, 1.1) {Case 2: $T/T_0$ $\to$ explicit};
\node[inp] (i2) at (0, 2.2) {Case 3: $P/P_0$ $\to$ explicit};
\node[inp] (i3) at (0, 3.3) {Case 4: $\rho/\rho_0$ $\to$ explicit};
\node[inp] (i4) at (0, 4.4) {Case 5: $A/A^*$ $\to$ N--R+toggle};
\node[inp] (i5) at (0, 5.5) {Case 6: $\mu$ $\to$ explicit};
\node[inp] (i6) at (0, 6.6) {Case 7: $\nu$ $\to$ Newton--Raphson};
\node[note] (tog) at (4.8, 4.4) {Sub/Sup\\toggle};
%--- Engine ---
\node[eng, minimum height=7.2cm] (eng) at (7.6, 3.3)
    {\textbf{IsentropicEngine}\\[3pt]
    \texttt{fromMach(M,\,$\gamma$)}};
%--- Outputs ---
\node[outn] (o0) at (12.2, 0.0) {$M$};
\node[outn] (o1) at (12.2, 1.1) {$T/T_0$};
\node[outn] (o2) at (12.2, 2.2) {$P/P_0$};
\node[outn] (o3) at (12.2, 3.3) {$\rho/\rho_0$};
\node[outn] (o4) at (12.2, 4.4) {$A/A^*$};
\node[outn] (o5) at (12.2, 5.5) {$\mu$};
\node[outn] (o6) at (12.2, 6.6) {$nu$};
%--- Input arrows ---
\draw[arr] (i0.east) -- (eng.west |- i0);
\draw[arr] (i1.east) -- (eng.west |- i1);
\draw[arr] (i2.east) -- (eng.west |- i2);
\draw[arr] (i3.east) -- (eng.west);
\draw[arr] (i4.east) -- (eng.west |- i4);
\draw[arr] (i5.east) -- (eng.west |- i5);
\draw[arr] (i6.east) -- (eng.west |- i6);
\draw[darr] (tog.east) -- (i4.east |- tog) -- (i4.east);
%--- Bus + output arrows ---
\coordinate (bx) at ($(eng.east)+(0.5,0)$);
\draw[bus] (eng.east) -- (bx);
\draw[bus] (bx |- o0) -- (bx |- o6);
\foreach \n in {o0,o1,o2,o3,o4,o5,o6} {\draw[arr] (bx |- \n) -- (\n.west);}
\end{tikzpicture}
\end{center}

\noindent Any one input drives all outputs. Explicit inputs ($T/T_0$, $P/P_0$, $\rho/\rho_0$, $\mu$) are inverted analytically. $A/A^*$ uses Newton--Raphson iteration with sub/supersonic toggle. $\nu$ uses Newton--Raphson with analytic derivative. Note: $\mu$ and $\nu$ are only defined for $M\ge1$.

